% !TeX program = xelatex
% Encoding: UTF-8
\documentclass[12pt,a4paper]{article}

\usepackage{fontspec}
\usepackage{xeCJK}
\usepackage{amsmath,amssymb,amsthm,mathtools}
\usepackage{geometry}
\usepackage{enumitem}
\usepackage{hyperref}

\geometry{margin=2.0cm}
\setmainfont{Latin Modern Roman}
\setCJKmainfont{標楷體}

\setlength{\parindent}{2em}
\setlength{\parskip}{0.5em}

% Make displayed matrices easier to read
\renewcommand{\arraystretch}{1.35}
\setlength{\arraycolsep}{9pt}

\newtheorem{theorem}{定理}

\title{\textbf{Schur's Theorem（Schur 定理）}\\
\large 歸納法證明}
\author{}
\date{}

\begin{document}
\maketitle

\section{酉矩陣}

對一個 $n\times n$ 的複數矩陣 $Q$，若
\[
    \overline{Q}^{\,T}=Q^{-1},
\]
則稱 $Q$ 為\textbf{酉矩陣（unitary matrix）}。

如果一個酉矩陣同時也是實矩陣，那麼
\[
    \overline{Q}=Q,
\]
因此
\[
    Q^T=Q^{-1}.
\]
所以實數情形下的酉矩陣就是\textbf{正交矩陣（orthogonal matrix）}。

\section{Schur 定理}

\begin{theorem}[Schur's Theorem]
    每一個方陣都可以分解成
    \[
        \boxed{
            A=QTQ^{-1}=QT\overline{Q}^{\,T},
        }
    \]
    其中 $T$ 是上三角矩陣，而 $Q$ 是酉矩陣。

    如果 $A$ 的所有特徵值皆為實數，則 $Q$ 與 $T$ 都可以選成實矩陣。
    此時
    \[
        Q^T=Q^{-1},
    \]
    也就是說 $Q$ 為正交矩陣。
\end{theorem}

\section{證明}

我們使用數學歸納法證明。

\subsection*{1. 基礎情形}

當 $n=1$ 時，結論顯然成立：
\[
    a=1\cdot a\cdot 1^{-1}.
\]

\subsection*{2. 歸納假設}

假設 Schur 定理對所有 $k\times k$ 矩陣成立，
現在令 $A$ 為一個 $(k+1)\times(k+1)$ 矩陣。

令 $\lambda_1$ 為 $A$ 的一個特徵值，
並令 $q_1$ 為對應的單位特徵向量，因此
\[
    Aq_1=\lambda_1q_1,
    \qquad
    \|q_1\|=1.
\]

利用 Gram--Schmidt 過程，可以找到
\[
    q_2,q_3,\ldots,q_{k+1},
\]
使得
\[
    q_1,q_2,\ldots,q_{k+1}
\]
形成 $\mathbb{C}^{k+1}$ 的一組正交規範基。

令
\[
    Q_1=
    \left[
        \begin{array}{cccc}
            \;q_1\; & \;q_2\; & \cdots & \;q_{k+1}\;
        \end{array}
        \right].
\]

因為 $q_1,q_2,\ldots,q_{k+1}$ 是一組正交規範基，
所以 $Q_1$ 是酉矩陣。

\subsection*{3. 計算 $\overline{Q}_1^{\,T}AQ_1$}

先把矩陣乘法完整寫開：
\[
    \overline{Q}_1^{\,T}AQ_1
    =
    \left[
        \begin{array}{c}
            \overline{q}_1^{\,T} \\[3pt]
            \overline{q}_2^{\,T} \\[3pt]
            \vdots               \\[3pt]
            \overline{q}_{k+1}^{\,T}
        \end{array}
        \right]
    \left[
        \begin{array}{cccc}
            Aq_1 & Aq_2 & \cdots & Aq_{k+1}
        \end{array}
        \right].
\]

因為
\[
    Aq_1=\lambda_1q_1,
\]
所以
\[
    \overline{Q}_1^{\,T}AQ_1
    =
    \left[
        \begin{array}{c}
            \overline{q}_1^{\,T} \\[3pt]
            \overline{q}_2^{\,T} \\[3pt]
            \vdots               \\[3pt]
            \overline{q}_{k+1}^{\,T}
        \end{array}
        \right]
    \left[
        \begin{array}{cccc}
            \lambda_1q_1 & Aq_2 & \cdots & Aq_{k+1}
        \end{array}
        \right].
\]

因此可以直接展開成
\[
    \overline{Q}_1^{\,T}AQ_1
    =
    \left[
        \begin{array}{cccc}
            \overline{q}_1^{\,T}\lambda_1q_1
                   &
            \overline{q}_1^{\,T}Aq_2
                   &
            \cdots
                   &
            \overline{q}_1^{\,T}Aq_{k+1}
            \\[5pt]
            \overline{q}_2^{\,T}\lambda_1q_1
                   &
            \overline{q}_2^{\,T}Aq_2
                   &
            \cdots
                   &
            \overline{q}_2^{\,T}Aq_{k+1}
            \\[5pt]
            \vdots & \vdots & \ddots & \vdots
            \\[5pt]
            \overline{q}_{k+1}^{\,T}\lambda_1q_1
                   &
            \overline{q}_{k+1}^{\,T}Aq_2
                   &
            \cdots
                   &
            \overline{q}_{k+1}^{\,T}Aq_{k+1}
        \end{array}
        \right].
\]

由正交規範性，
\[
    \overline{q}_1^{\,T}q_1=1,
    \qquad
    \overline{q}_j^{\,T}q_1=0
    \quad(j=2,\ldots,k+1),
\]
因此第一欄化為
\[
    \left[
        \begin{array}{c}
            \lambda_1 \\[4pt]
            0         \\[4pt]
            \vdots    \\[4pt]
            0
        \end{array}
        \right].
\]

所以整個矩陣可寫成
\[
    \boxed{
        \overline{Q}_1^{\,T}AQ_1
        =
        \left[
            \begin{array}{ccccc}
                \lambda_1
                 &
                *
                 &
                *
                 &
                \cdots
                 &
                *
                \\[5pt]
                0
                 &
                 &
                 &
                 &
                \\[-3pt]
                0
                 &
                 &
                A_2
                 &
                 &
                \\[-3pt]
                \vdots
                 &
                 &
                 &
                 &
                \\[-3pt]
                0
                 &
                 &
                 &
                 &
            \end{array}
            \right]
    }
\]
其中 $A_2$ 是一個 $k\times k$ 矩陣。

這裡最重要的是：
\[
    \boxed{\text{第一欄除了 }\lambda_1\text{ 以外，其餘元素全部為 }0.}
\]

\subsection*{4. 對 $A_2$ 使用歸納假設}

因為 $A_2$ 是一個 $k\times k$ 矩陣，
依照歸納假設，存在酉矩陣 $Q_2$ 與上三角矩陣 $T_2$，使得
\[
    A_2=Q_2T_2\overline{Q}_2^{\,T},
\]
因此
\[
    \overline{Q}_2^{\,T}A_2Q_2=T_2.
\]

令
\[
    Q
    =
    Q_1
    \left[
        \begin{array}{ccccc}
            1      & 0 & 0   & \cdots & 0 \\[4pt]
            0      &   &     &        &   \\[-3pt]
            0      &   & Q_2 &        &   \\[-3pt]
            \vdots &   &     &        &   \\[-3pt]
            0      &   &     &        &
        \end{array}
        \right].
\]

為方便表示，設
\[
    D=
    \left[
        \begin{array}{ccccc}
            1      & 0 & 0   & \cdots & 0 \\[4pt]
            0      &   &     &        &   \\[-3pt]
            0      &   & Q_2 &        &   \\[-3pt]
            \vdots &   &     &        &   \\[-3pt]
            0      &   &     &        &
        \end{array}
        \right],
    \qquad
    Q=Q_1D.
\]

\subsection*{5. 證明 $Q$ 是酉矩陣}

因為 $Q_2$ 是酉矩陣，
\[
    Q_2\overline{Q}_2^{\,T}=I_k.
\]

因此
\[
    D\overline{D}^{\,T}
    =
    \left[
        \begin{array}{ccccc}
            1      & 0 & 0   & \cdots & 0 \\
            0      &   &     &        &   \\
            0      &   & Q_2 &        &   \\
            \vdots &   &     &        &   \\
            0      &   &     &        &
        \end{array}
        \right]
    \left[
        \begin{array}{ccccc}
            1      & 0 & 0                    & \cdots & 0 \\
            0      &   &                      &        &   \\
            0      &   & \overline{Q}_2^{\,T} &        &   \\
            \vdots &   &                      &        &   \\
            0      &   &                      &        &
        \end{array}
        \right]
    =
    \left[
        \begin{array}{ccccc}
            1      & 0 & 0   & \cdots & 0 \\
            0      &   &     &        &   \\
            0      &   & I_k &        &   \\
            \vdots &   &     &        &   \\
            0      &   &     &        &
        \end{array}
        \right]
    =
    I_{k+1}.
\]

又因為 $Q_1$ 是酉矩陣，
\[
    Q_1\overline{Q}_1^{\,T}=I_{k+1}.
\]

故
\[
    \begin{aligned}
        Q\overline{Q}^{\,T}
         & =
        Q_1D\overline{D}^{\,T}\overline{Q}_1^{\,T} \\
         & =
        Q_1I_{k+1}\overline{Q}_1^{\,T}             \\
         & =
        Q_1\overline{Q}_1^{\,T}                    \\
         & =
        I_{k+1}.
    \end{aligned}
\]

因此 $Q$ 是酉矩陣。

\subsection*{6. 計算 $\overline{Q}^{\,T}AQ$}

由
\[
    Q=Q_1D,
\]
可得
\[
    \overline{Q}^{\,T}AQ
    =
    \overline{D}^{\,T}\,
    \overline{Q}_1^{\,T}AQ_1\,
    D.
\]

完整寫成矩陣：
\[
    \overline{Q}^{\,T}AQ
    =
    \left[
        \begin{array}{ccccc}
            1      & 0 & 0                    & \cdots & 0 \\
            0      &   &                      &        &   \\
            0      &   & \overline{Q}_2^{\,T} &        &   \\
            \vdots &   &                      &        &   \\
            0      &   &                      &        &
        \end{array}
        \right]
    \left[
        \begin{array}{ccccc}
            \lambda_1 & * & *   & \cdots & * \\
            0         &   &     &        &   \\
            0         &   & A_2 &        &   \\
            \vdots    &   &     &        &   \\
            0         &   &     &        &
        \end{array}
        \right]
    \left[
        \begin{array}{ccccc}
            1      & 0 & 0   & \cdots & 0 \\
            0      &   &     &        &   \\
            0      &   & Q_2 &        &   \\
            \vdots &   &     &        &   \\
            0      &   &     &        &
        \end{array}
        \right].
\]

因此
\[
    \overline{Q}^{\,T}AQ
    =
    \left[
        \begin{array}{ccccc}
            \lambda_1 & * & *                          & \cdots & * \\
            0         &   &                            &        &   \\
            0         &   & \overline{Q}_2^{\,T}A_2Q_2 &        &   \\
            \vdots    &   &                            &        &   \\
            0         &   &                            &        &
        \end{array}
        \right].
\]

再利用
\[
    \overline{Q}_2^{\,T}A_2Q_2=T_2,
\]
得到
\[
    \boxed{
        \overline{Q}^{\,T}AQ
        =
        \left[
            \begin{array}{ccccc}
                \lambda_1 & * & *   & \cdots & * \\
                0         &   &     &        &   \\
                0         &   & T_2 &        &   \\
                \vdots    &   &     &        &   \\
                0         &   &     &        &
            \end{array}
            \right]
        =T.
    }
\]

因為 $T_2$ 是上三角矩陣，所以 $T$ 也是上三角矩陣。

因此
\[
    \overline{Q}^{\,T}AQ=T.
\]

又因為
\[
    \overline{Q}^{\,T}=Q^{-1},
\]
所以
\[
    Q^{-1}AQ=T.
\]

左乘 $Q$、右乘 $Q^{-1}$，得到
\[
    \boxed{
        A=QTQ^{-1}=QT\overline{Q}^{\,T}.
    }
\]

因此 Schur 定理對 $(k+1)\times(k+1)$ 矩陣成立。

由數學歸納法可知，Schur 定理對所有 $n\times n$ 複數矩陣皆成立。

\end{document}
